\documentclass[UTF8, 12pt, fontset=fandol, oneside]{ctexart}
\usepackage{researchnotes}

\title{强化学习：动态规划}
\author{}
\date{}

\begin{document}

\maketitle

\section*{第 4 章\quad 动态规划}

动态规划（dynamic programming, DP）这个术语指的是一组算法；在给定环境作为马尔可夫决策过程（Markov decision process, MDP）的完美模型时，这些算法可以用来计算最优策略。
经典 DP 算法在强化学习中的实用性有限，这既是因为它们假设存在完美模型，也是因为它们的计算代价很高；但它们在理论上仍然很重要。
DP 为理解本书其余部分所介绍的方法提供了必要基础。
事实上，所有这些方法都可以看作是在试图达到与 DP 大致相同的效果，只是计算量更少，而且不假设环境存在完美模型。

我们通常假设环境是有限 MDP。
也就是说，我们假设它的状态集合、动作集合和奖励集合 $\mathcal{S}$、$\mathcal{A}$ 与 $\mathcal{R}$ 都是有限的，并且它的动态由一组概率 $p(s',r|s,a)$ 给出，其中 $s\in\mathcal{S}$、$a\in\mathcal{A}(s)$、$r\in\mathcal{R}$ 且 $s'\in\mathcal{S}^+$（$\mathcal{S}^+$ 是 $\mathcal{S}$ 加上一个终止状态，如果问题是情节式的）。
虽然 DP 的思想可以应用于具有连续状态空间和动作空间的问题，但精确解只在一些特殊情形中才可能得到。
对于具有连续状态和动作的任务，一种常见的近似求解方式是先把状态空间和动作空间量化，然后应用有限状态的 DP 方法。
我们将在第 9 章考察的方法可以用于连续问题，并且是对这种方法的重要扩展。

DP 的关键思想，也是强化学习一般意义上的关键思想，是使用价值函数来组织和构造对好策略的搜索。
在本章中，我们说明如何用 DP 来计算第 3 章中定义的价值函数。
如那里所讨论的，一旦我们找到了最优价值函数 $v_*$ 或 $q_*$，就可以很容易地得到最优策略；它们满足贝尔曼最优方程：
\begin{equation}
\begin{aligned}
v_*(s)
&= \max_a \mathbb{E}\!\left[R_{t+1}+\gamma v_*(S_{t+1}) \mid S_t=s,A_t=a\right] \\
&= \max_a \sum_{s',r} p(s',r|s,a)\left[r+\gamma v_*(s')\right],
\end{aligned}
\tag{4.1}
\end{equation}
或者
\begin{equation}
\begin{aligned}
q_*(s,a)
&= \mathbb{E}\!\left[R_{t+1}+\gamma\max_{a'}q_*(S_{t+1},a') \mid S_t=s,A_t=a\right] \\
&= \sum_{s',r} p(s',r|s,a)\left[r+\gamma\max_{a'}q_*(s',a')\right],
\end{aligned}
\tag{4.2}
\end{equation}
对所有 $s\in\mathcal{S}$、$a\in\mathcal{A}(s)$ 和 $s'\in\mathcal{S}^+$ 成立。
正如我们将看到的，DP 算法是通过把这样的贝尔曼方程转化为赋值语句得到的，也就是转化为用于改进所需价值函数近似的更新规则。

\subsection*{4.1\quad 策略评估（预测）}

首先，我们考虑如何为任意策略 $\pi$ 计算状态价值函数 $v_\pi$。
在 DP 文献中，这称为\textbf{策略评估}（policy evaluation）。
我们也把它称为\textbf{预测问题}（prediction problem）。
回忆第 3 章，对所有 $s\in\mathcal{S}$，
\begin{align}
v_\pi(s)
&\doteq \mathbb{E}_\pi[G_t\mid S_t=s] \\
&= \mathbb{E}_\pi[R_{t+1}+\gamma G_{t+1}\mid S_t=s] \qquad \text{（由 (3.9)）}\\
&= \mathbb{E}_\pi[R_{t+1}+\gamma v_\pi(S_{t+1})\mid S_t=s] \tag{4.3}\\
&= \sum_a \pi(a|s)\sum_{s',r}p(s',r|s,a)\left[r+\gamma v_\pi(s')\right].
\tag{4.4}
\end{align}
这里，$\pi(a|s)$ 是在策略 $\pi$ 下、状态 $s$ 中采取动作 $a$ 的概率；期望的下标 $\pi$ 表明这些期望是在遵循 $\pi$ 的条件下取的。
只要 $\gamma<1$，或者策略 $\pi$ 保证从所有状态最终终止，$v_\pi$ 的存在性与唯一性就能得到保证。

如果环境的动态完全已知，那么 (4.4) 就是一个由 $|\mathcal{S}|$ 个未知量（即 $v_\pi(s),s\in\mathcal{S}$）组成的 $|\mathcal{S}|$ 个联立线性方程系统。
原则上，求解它是一项直接但繁琐的计算。
就我们的目的而言，迭代求解方法最为合适。
考虑一列近似价值函数 $v_0,v_1,v_2,\ldots$，其中每个函数都把 $\mathcal{S}^+$ 映射到 $\mathbb{R}$（实数）。
初始近似 $v_0$ 可以任意选择（但终止状态若存在，则必须赋值为 $0$），每一个后继近似都通过把 $v_\pi$ 的贝尔曼方程 (4.4) 用作更新规则来得到：
\begin{equation}
\begin{aligned}
v_{k+1}(s)
&\doteq \mathbb{E}_\pi[R_{t+1}+\gamma v_k(S_{t+1})\mid S_t=s] \\
&= \sum_a \pi(a|s)\sum_{s',r}p(s',r|s,a)\left[r+\gamma v_k(s')\right],
\end{aligned}
\tag{4.5}
\end{equation}
对所有 $s\in\mathcal{S}$ 成立。
显然，$v_k=v_\pi$ 是这个更新规则的一个不动点，因为 $v_\pi$ 的贝尔曼方程保证了此时的等式成立。
事实上，在保证 $v_\pi$ 存在的相同条件下，可以一般性地证明序列 $\{v_k\}$ 会随着 $k\to\infty$ 收敛到 $v_\pi$。
这个算法称为\textbf{迭代策略评估}。

为了由 $v_k$ 产生每一个后继近似 $v_{k+1}$，迭代策略评估对每个状态 $s$ 应用同一个操作：它用一个新值替换 $s$ 的旧值；这个新值由 $s$ 的后继状态的旧值、期望即时奖励以及被评估策略下所有可能的一步转移得到。
我们把这类操作称为\textbf{期望更新}（expected update）。
迭代策略评估的每一次迭代都对每个状态的价值更新一次，从而产生新的近似价值函数 $v_{k+1}$。
根据被更新的是一个状态（如这里）还是一个状态--动作对，以及后继状态估计值的精确组合方式，期望更新有若干不同类型。
DP 算法中的所有更新都称为期望更新，因为它们基于对所有可能下一状态的期望，而不是基于一个采样得到的下一状态。
更新的性质可以像上面那样用方程表达，也可以像第 3 章中介绍的那样用备份图表达。
例如，与迭代策略评估中所用期望更新相对应的备份图见第 59 页。

若要按照 (4.5) 写一个顺序计算机程序来实现迭代策略评估，就需要使用两个数组：一个保存旧值 $v_k(s)$，另一个保存新值 $v_{k+1}(s)$。
有了两个数组，新值就可以逐个由旧值计算出来，而旧值不会被改变。
当然，使用一个数组并“原地”更新这些值更容易，也就是每个新值立即覆盖旧值。
于是，根据状态更新的顺序不同，(4.5) 右端有时会使用新值而不是旧值。
这个原地算法也会收敛到 $v_\pi$；事实上，通常它比双数组版本收敛得更快，这也合乎预期，因为它一有新数据就立即使用。
我们把这些更新看作是在状态空间中进行一次\textbf{扫掠}（sweep）。
对原地算法来说，状态在扫掠中被更新的顺序会显著影响收敛速度。
当我们思考 DP 算法时，通常心里想的是原地版本。

下面的框中给出了迭代策略评估的一个完整原地版本伪代码。
注意它如何处理终止。
形式上，迭代策略评估只有在极限中才收敛，但实践中必须在此之前停止。
该伪代码在每一次扫掠后检查量 $\max_{s\in\mathcal{S}}|v_{k+1}(s)-v_k(s)|$，并在它足够小时停止。

\begin{center}
\fbox{\begin{minipage}{0.92\textwidth}
\noindent\textbf{迭代策略评估}，用于估计 $V\approx v_\pi$

\medskip
\noindent 输入 $\pi$，即待评估的策略。

\noindent 算法参数：一个小阈值 $\theta>0$，用于决定估计精度。

\noindent 对所有 $s\in\mathcal{S}^+$ 初始化 $V(s)$，可任意取值，但 $V(\mathrm{terminal})=0$。

\medskip
\noindent 循环：

\quad $\Delta \leftarrow 0$

\quad 对每个 $s\in\mathcal{S}$ 循环：

\qquad $v\leftarrow V(s)$

\qquad $V(s)\leftarrow \sum_a\pi(a|s)\sum_{s',r}p(s',r|s,a)\left[r+\gamma V(s')\right]$

\qquad $\Delta\leftarrow \max(\Delta,|v-V(s)|)$

\quad 直到 $\Delta<\theta$
\end{minipage}}
\end{center}

\noindent\textbf{例 4.1}
考虑下图所示的 $4\times4$ 网格世界。

\begin{figure}[htbp]
\centering
\includegraphics[width=0.82\textwidth]{tex/强化学习/figures/rl5-example-4-1-grid.png}
\caption*{例 4.1 的 $4\times4$ 网格世界示意图。}
\end{figure}

非终止状态为 $\mathcal{S}=\{1,2,\ldots,14\}$。
每个状态中有四个可能动作，
\[
\mathcal{A}=\{\mathrm{up},\mathrm{down},\mathrm{right},\mathrm{left}\},
\]
它们确定性地引起相应的状态转移；但如果某个动作会把智能体带出网格，则实际上让状态保持不变。
因此，例如 $p(6,-1|5,\mathrm{right})=1$，$p(7,-1|7,\mathrm{right})=1$，并且对所有 $r\in\mathcal{R}$ 有 $p(10,r|5,\mathrm{right})=0$。
这是一个无折扣的情节式任务。
在达到终止状态之前，所有转移的奖励都是 $-1$。
终止状态在图中以阴影表示（虽然它在两个位置显示，但形式上它是同一个状态）。
因此，对所有状态 $s,s'$ 和动作 $a$，期望奖励函数都是 $r(s,a,s')=-1$。
假设智能体遵循等概率随机策略（所有动作等可能）。
图 4.1 的左侧显示了由迭代策略评估计算得到的价值函数序列 $\{v_k\}$。
最终估计事实上就是 $v_\pi$；在这个例子中，它给出了从每个状态到终止状态的期望步数的相反数。

\noindent\textbf{习题 4.1}
在例 4.1 中，如果 $\pi$ 是等概率随机策略，那么 $q_\pi(11,\mathrm{down})$ 是多少？
$q_\pi(7,\mathrm{down})$ 是多少？

\noindent\textbf{习题 4.2}
在例 4.1 中，假设在状态 13 的正下方增加一个新状态 15，并且它的动作 $\mathrm{left}$、$\mathrm{up}$、$\mathrm{right}$ 和 $\mathrm{down}$ 分别把智能体带到状态 12、13、14 和 15。
假设来自原有状态的转移保持不变。
那么，在等概率随机策略下，$v_\pi(15)$ 是多少？
现在再假设状态 13 的动态也发生改变，使得从状态 13 执行动作 $\mathrm{down}$ 会把智能体带到新状态 15。
在这种情况下，等概率随机策略下的 $v_\pi(15)$ 是多少？

\noindent\textbf{习题 4.3}
对于动作价值函数 $q_\pi$ 以及由序列 $q_0,q_1,q_2,\ldots$ 给出的逐次近似，与 (4.3)、(4.4) 和 (4.5) 类似的方程是什么？

\begin{figure}[htbp]
\centering
\includegraphics[width=0.62\textwidth]{tex/强化学习/figures/rl5-fig-4-1.png}
\caption{小网格世界上迭代策略评估的收敛。左列是随机策略（所有动作等可能）的状态价值函数近似序列。右列是与这些价值函数估计相对应的贪心策略序列（所有达到最大值的动作都用箭头表示，数字四舍五入到两位有效数字）。最后一个策略只保证优于随机策略，但在这个例子中，它以及第三次迭代之后的所有策略都是最优的。}
\end{figure}

\subsection*{4.2\quad 策略改进}

我们计算一个策略的价值函数，是为了帮助找到更好的策略。
假设我们已经确定了某个任意确定性策略 $\pi$ 的价值函数 $v_\pi$。
对于某个状态 $s$，我们想知道是否应该改变策略，使其确定性地选择某个动作 $a\ne\pi(s)$。
我们知道从 $s$ 开始遵循当前策略有多好，即 $v_\pi(s)$；但改用新策略会更好还是更差呢？
回答这个问题的一种方式是考虑在 $s$ 中选择 $a$，之后继续遵循已有策略 $\pi$。
这种行为方式的价值为
\begin{equation}
\begin{aligned}
q_\pi(s,a)
&\doteq \mathbb{E}\!\left[R_{t+1}+\gamma v_\pi(S_{t+1}) \mid S_t=s,A_t=a\right] \\
&= \sum_{s',r}p(s',r|s,a)\left[r+\gamma v_\pi(s')\right].
\end{aligned}
\tag{4.6}
\end{equation}

关键判据是它大于还是小于 $v_\pi(s)$。
如果它更大，也就是说，如果在 $s$ 中先选择一次 $a$、之后遵循 $\pi$ 比始终遵循 $\pi$ 更好，那么我们会预期每次遇到 $s$ 都选择 $a$ 仍然会更好，而且新策略整体上实际上会更好。

这一点成立，是一个称为\textbf{策略改进定理}的一般结果的特殊情形。
令 $\pi$ 和 $\pi'$ 是任意一对确定性策略，并且对所有 $s\in\mathcal{S}$ 都有
\begin{equation}
q_\pi(s,\pi'(s))\ge v_\pi(s).
\tag{4.7}
\end{equation}
那么策略 $\pi'$ 必须至少与 $\pi$ 一样好，或者更好。
也就是说，它必须从所有状态 $s\in\mathcal{S}$ 得到不小于 $\pi$ 的期望回报：
\begin{equation}
v_{\pi'}(s)\ge v_\pi(s).
\tag{4.8}
\end{equation}
此外，如果 (4.7) 在任一状态处是严格不等式，那么 (4.8) 在该状态处也必然是严格不等式。
这个结果特别适用于上一段中考虑的两个策略：原确定性策略 $\pi$，以及一个改变后的策略 $\pi'$；后者除了 $\pi'(s)=a\ne\pi(s)$ 之外与 $\pi$ 相同。
显然，在除 $s$ 之外的所有状态上 (4.7) 都成立。
因此，如果 $q_\pi(s,a)>v_\pi(s)$，那么改变后的策略确实优于 $\pi$。

策略改进定理的证明思想很容易理解。
从 (4.7) 出发，我们不断用 (4.6) 展开 $q_\pi$ 一侧，并反复应用 (4.7)，直到得到 $v_{\pi'}(s)$：
\[
\begin{aligned}
v_\pi(s)
&\le q_\pi(s,\pi'(s))\\
&= \mathbb{E}[R_{t+1}+\gamma v_\pi(S_{t+1})\mid S_t=s,A_t=\pi'(s)] \qquad \text{（由 (4.6)）}\\
&= \mathbb{E}_{\pi'}[R_{t+1}+\gamma v_\pi(S_{t+1})\mid S_t=s]\\
&\le \mathbb{E}_{\pi'}[R_{t+1}+\gamma q_\pi(S_{t+1},\pi'(S_{t+1}))\mid S_t=s] \qquad \text{（由 (4.7)）}\\
&= \mathbb{E}_{\pi'}[R_{t+1}+\gamma\mathbb{E}_{\pi'}[R_{t+2}+\gamma v_\pi(S_{t+2})\mid S_{t+1},A_{t+1}=\pi'(S_{t+1})]\mid S_t=s]\\
&= \mathbb{E}_{\pi'}[R_{t+1}+\gamma R_{t+2}+\gamma^2v_\pi(S_{t+2})\mid S_t=s]\\
&\le \mathbb{E}_{\pi'}[R_{t+1}+\gamma R_{t+2}+\gamma^2R_{t+3}+\gamma^3v_\pi(S_{t+3})\mid S_t=s]\\
&\quad\vdots\\
&\le \mathbb{E}_{\pi'}[R_{t+1}+\gamma R_{t+2}+\gamma^2R_{t+3}+\gamma^3R_{t+4}+\cdots\mid S_t=s]\\
&= v_{\pi'}(s).
\end{aligned}
\]

到目前为止，我们已经看到，在给定一个策略及其价值函数时，可以很容易地评估在单个状态把策略改为某个特定动作的效果。
自然的扩展是考虑在所有状态、对所有可能动作都进行改变，在每个状态选择按照 $q_\pi(s,a)$ 看起来最好的动作。
换句话说，考虑由下式给出的新贪心策略 $\pi'$：
\begin{equation}
\begin{aligned}
\pi'(s)
&\doteq \operatorname*{arg\,max}_a q_\pi(s,a)\\
&= \operatorname*{arg\,max}_a \mathbb{E}[R_{t+1}+\gamma v_\pi(S_{t+1})\mid S_t=s,A_t=a] \\
&= \operatorname*{arg\,max}_a \sum_{s',r}p(s',r|s,a)\left[r+\gamma v_\pi(s')\right].
\end{aligned}
\tag{4.9}
\end{equation}
这里，$\operatorname*{arg\,max}_a$ 表示使其后表达式达到最大值的 $a$ 的取值（若有并列，则任意打破并列）。
贪心策略采取的是按 $v_\pi$ 看来经过一步前瞻后短期最好的动作。
按照构造，贪心策略满足策略改进定理 (4.7) 的条件，因此我们知道它至少与原策略一样好，或者更好。
用原策略的价值函数把新策略改成相对于它的贪心策略，从而产生一个优于原策略的新策略，这个过程称为\textbf{策略改进}。

假设新的贪心策略 $\pi'$ 与旧策略 $\pi$ 一样好，但并不更好。
于是 $v_\pi=v_{\pi'}$，并且由 (4.9) 可知，对所有 $s\in\mathcal{S}$，
\[
\begin{aligned}
v_{\pi'}(s)
&= \max_a \mathbb{E}[R_{t+1}+\gamma v_{\pi'}(S_{t+1})\mid S_t=s,A_t=a]\\
&= \max_a \sum_{s',r}p(s',r|s,a)\left[r+\gamma v_{\pi'}(s')\right].
\end{aligned}
\]
但这正是贝尔曼最优方程 (4.1)，因此 $v_{\pi'}$ 必须是 $v_*$，并且 $\pi$ 与 $\pi'$ 都必须是最优策略。
因此，除非原策略已经是最优的，否则策略改进必定给出一个严格更好的策略。

到目前为止，本节只考虑了确定性策略这一特殊情形。
一般情况下，随机策略 $\pi$ 为每个状态 $s$ 中采取每个动作 $a$ 指定概率 $\pi(a|s)$。
我们不展开细节，但本节的所有思想实际上都很容易扩展到随机策略。
特别是，策略改进定理按上述形式也适用于随机情形。
另外，如果在诸如 (4.9) 的策略改进步骤中出现并列，也就是说，如果有若干动作同时达到最大值，那么在随机情形中我们不必从它们中选出单个动作。
相反，可以把新贪心策略中被选择的概率分配给每个达到最大值的动作。
只要所有次优动作得到的概率为零，任何分配方案都是允许的。

图 4.1 的最后一行展示了随机策略的策略改进示例。
这里原策略 $\pi$ 是等概率随机策略，而新策略 $\pi'$ 相对于 $v_\pi$ 是贪心的。
价值函数 $v_\pi$ 显示在左下图，可能的 $\pi'$ 集合显示在右下图。
在 $\pi'$ 图中带有多个箭头的状态，就是 (4.9) 中有多个动作达到最大值的状态；允许在这些动作之间以任意方式分配概率。
通过观察可以看出，任何这种策略的价值函数 $v_{\pi'}(s)$ 在所有 $s\in\mathcal{S}$ 上都只能是 $-1$、$-2$ 或 $-3$，而 $v_\pi(s)$ 至多为 $-14$。
因此，对所有 $s\in\mathcal{S}$，都有 $v_{\pi'}(s)\ge v_\pi(s)$，这说明了策略改进。
虽然在这个例子中新策略 $\pi'$ 恰好是最优的，但一般来说只能保证得到改进。

\subsection*{4.3\quad 策略迭代}

一旦使用 $v_\pi$ 改进了策略 $\pi$，得到更好的策略 $\pi'$，我们就可以计算 $v_{\pi'}$，并再次改进它以得到更好的策略 $\pi''$。
因此，我们可以得到一列单调改进的策略和价值函数：
\[
\pi_0 \xrightarrow{\mathrm{E}} v_{\pi_0} \xrightarrow{\mathrm{I}} \pi_1
\xrightarrow{\mathrm{E}} v_{\pi_1} \xrightarrow{\mathrm{I}} \pi_2
\xrightarrow{\mathrm{E}} \cdots \xrightarrow{\mathrm{I}} \pi_*
\xrightarrow{\mathrm{E}} v_* ,
\]
其中 $\xrightarrow{\mathrm{E}}$ 表示策略评估，$\xrightarrow{\mathrm{I}}$ 表示策略改进。
每个策略都保证是对前一个策略的严格改进（除非它已经是最优的）。
由于有限 MDP 只有有限多个策略，这个过程必定在有限次迭代内收敛到一个最优策略和最优价值函数。

这种寻找最优策略的方法称为\textbf{策略迭代}。
下面的框中给出了完整算法。
注意，每一次策略评估本身也是一次迭代计算，并且从前一个策略的价值函数开始。
这通常会极大提高策略评估的收敛速度（大概是因为从一个策略到下一个策略，价值函数变化不大）。

\begin{center}
\fbox{\begin{minipage}{0.92\textwidth}
\noindent\textbf{策略迭代（使用迭代策略评估）}，用于估计 $\pi\approx\pi_*$

\medskip
\noindent 1. 初始化

\quad 对所有 $s\in\mathcal{S}$，任意取 $V(s)\in\mathbb{R}$ 且 $\pi(s)\in\mathcal{A}(s)$。

\medskip
\noindent 2. 策略评估

\quad 循环：

\qquad $\Delta\leftarrow 0$

\qquad 对每个 $s\in\mathcal{S}$ 循环：

\qquad\quad $v\leftarrow V(s)$

\qquad\quad $V(s)\leftarrow \sum_{s',r}p(s',r|s,\pi(s))\left[r+\gamma V(s')\right]$

\qquad\quad $\Delta\leftarrow \max(\Delta,|v-V(s)|)$

\quad 直到 $\Delta<\theta$（一个决定估计精度的小正数）。

\medskip
\noindent 3. 策略改进

\quad $\mathit{policy\mbox{-}stable}\leftarrow \mathit{true}$

\quad 对每个 $s\in\mathcal{S}$：

\qquad $\mathit{old\mbox{-}action}\leftarrow \pi(s)$

\qquad $\pi(s)\leftarrow \operatorname*{arg\,max}_a\sum_{s',r}p(s',r|s,a)\left[r+\gamma V(s')\right]$

\qquad 如果 $\mathit{old\mbox{-}action}\ne\pi(s)$，则 $\mathit{policy\mbox{-}stable}\leftarrow \mathit{false}$

\quad 如果 $\mathit{policy\mbox{-}stable}$，则停止并返回 $V\approx v_*$ 与 $\pi\approx\pi_*$；否则转到 2。
\end{minipage}}
\end{center}

\noindent\textbf{例 4.2：Jack 租车问题}
Jack 管理着一家全国性租车公司的两个网点。
每天，两个网点都会有若干顾客到来租车。
如果 Jack 有可用车辆，他就把车租出去，并从全国公司记入 $10$ 美元收益。
如果该网点没有车，那么这笔生意就损失了。
车辆在归还后的第二天可以再次出租。
为了确保车辆出现在需要它们的地方，Jack 可以在夜间在两个网点之间移动车辆，成本是每移动一辆车 $2$ 美元。
我们假设每个网点的租车请求数和还车数都是泊松随机变量，也就是说，该数等于 $n$ 的概率是 $\frac{\lambda^n}{n!}e^{-\lambda}$，其中 $\lambda$ 是期望数。
假设第一个和第二个网点的租车请求 $\lambda$ 分别为 $3$ 和 $4$，还车的 $\lambda$ 分别为 $3$ 和 $2$。
为了稍微简化问题，我们假设每个网点最多只能有 $20$ 辆车（多余的车会归还给全国公司，因此从问题中消失），并且每晚最多只能从一个网点向另一个网点移动五辆车。
我们取折扣率 $\gamma=0.9$，并把它表述为一个持续性有限 MDP：时间步为天，状态是一天结束时两个网点的车辆数，动作是夜间在两个网点之间移动的净车辆数。
图 4.2 显示了从从不移动车辆的策略开始，由策略迭代找到的策略序列。

\begin{figure}[htbp]
\centering
\includegraphics[width=0.92\textwidth]{tex/强化学习/figures/rl5-fig-4-2.png}
\caption{Jack 租车问题中由策略迭代找到的策略序列，以及最终状态价值函数。前五幅图显示了在每天结束时两个网点各有若干车辆的情况下，应从第一个网点向第二个网点移动的车辆数（负数表示从第二个网点转移到第一个网点）。每一个后继策略都是对前一策略的严格改进，最后一个策略是最优的。}
\end{figure}

策略迭代往往会以少得惊人的迭代次数收敛，Jack 租车的例子说明了这一点，图 4.1 的例子也说明了这一点。
图 4.1 的左下图给出了等概率随机策略的价值函数，右下图给出了该价值函数的一个贪心策略。
策略改进定理保证这些策略优于原随机策略。
不过在这个例子中，这些策略不仅更好，而且是最优的；它们沿最少步数路径前往终止状态。
在这个例子中，策略迭代只需要一次迭代就能找到最优策略。

\noindent\textbf{习题 4.4}
第 80 页的策略迭代算法有一个微妙的漏洞：如果策略不断在两个或更多个同样好的策略之间切换，它可能永远不会终止。
这作为教学用途没有问题，但不适合实际使用。
请修改伪代码，使得收敛得到保证。

\noindent\textbf{习题 4.5}
如何为动作价值定义策略迭代？
请给出一个用于计算 $q_*$ 的完整算法，类似第 80 页用于计算 $v_*$ 的算法。
请特别注意这道练习，因为其中涉及的思想会在本书其余部分反复使用。

\noindent\textbf{习题 4.6}
假设你被限制只能考虑 $\varepsilon$-soft 策略，也就是说，在每个状态 $s$ 中选择每个动作的概率至少为 $\varepsilon/|\mathcal{A}(s)|$。
请定性描述在第 80 页用于 $v_*$ 的策略迭代算法中，步骤 3、2 和 1 分别需要作出什么改变（按这个顺序）。

\noindent\textbf{习题 4.7（编程）}
为策略迭代写一个程序，并在下面这些改变下重新求解 Jack 租车问题。
Jack 在第一个网点的一名员工每天晚上乘公共汽车回家，并住在第二个网点附近。
她愿意免费把一辆车送到第二个网点。
每多移动一辆车仍然花费 $2$ 美元，反方向移动的所有车辆也都花费 $2$ 美元。
此外，Jack 在每个网点的停车空间有限。
如果某个网点夜间保留的车辆超过 $10$ 辆（在完成任何车辆移动之后），那么必须额外支付 $4$ 美元以使用第二个停车场（与那里保留多少辆车无关）。
这类非线性和任意动态经常出现在真实问题中，并且除了动态规划之外，其他优化方法通常很难处理。
为了检查你的程序，请先复现原问题给出的结果。

\subsection*{4.4\quad 价值迭代}

策略迭代的一个缺点是，它的每一次迭代都包含策略评估；而策略评估本身可能是一项漫长的迭代计算，需要对状态集合进行多次扫掠。
如果策略评估以迭代方式完成，那么精确收敛到 $v_\pi$ 只会在极限中发生。
我们必须等待精确收敛吗，还是可以提前停止？
图 4.1 的例子显然暗示，可以截断策略评估。
在那个例子中，前三次之后的策略评估迭代对相应的贪心策略没有影响。

事实上，在不失去策略迭代收敛保证的前提下，策略迭代中的策略评估步骤可以用几种方式截断。
一个重要的特殊情形是，策略评估只在一次扫掠（每个状态更新一次）后停止。
这个算法称为\textbf{价值迭代}。
它可以写成一个特别简单的更新操作，把策略改进和截断策略评估步骤结合在一起：
\begin{equation}
\begin{aligned}
v_{k+1}(s)
&\doteq \max_a \mathbb{E}\!\left[R_{t+1}+\gamma v_k(S_{t+1}) \mid S_t=s,A_t=a\right] \\
&= \max_a \sum_{s',r}p(s',r|s,a)\left[r+\gamma v_k(s')\right],
\end{aligned}
\tag{4.10}
\end{equation}
对所有 $s\in\mathcal{S}$ 成立。
对于任意 $v_0$，可以证明序列 $\{v_k\}$ 在保证 $v_*$ 存在的相同条件下收敛到 $v_*$。

理解价值迭代的另一种方式，是参考贝尔曼最优方程 (4.1)。
注意，价值迭代只是把贝尔曼最优方程转化成更新规则。
还要注意，价值迭代更新与策略评估更新 (4.5) 完全相同，只是它需要在所有动作上取最大值。
观察这种紧密关系的另一种方式，是比较这些算法的备份图：第 59 页是策略评估，图 3.4 左侧是价值迭代。
这二者分别是用于计算 $v_\pi$ 和 $v_*$ 的自然备份操作。

最后，让我们考虑价值迭代如何终止。
像策略评估一样，形式上价值迭代需要无限多次迭代才能精确收敛到 $v_*$。
在实践中，一旦价值函数在一次扫掠中的变化量很小，我们就停止。
下面的框给出了一个带有这种终止条件的完整算法。

\begin{center}
\fbox{\begin{minipage}{0.92\textwidth}
\noindent\textbf{价值迭代}，用于估计 $\pi\approx\pi_*$

\medskip
\noindent 算法参数：一个小阈值 $\theta>0$，用于决定估计精度。

\noindent 对所有 $s\in\mathcal{S}^+$ 初始化 $V(s)$，可任意取值，但 $V(\mathrm{terminal})=0$。

\medskip
\noindent 循环：

\quad $\Delta\leftarrow 0$

\quad 对每个 $s\in\mathcal{S}$ 循环：

\qquad $v\leftarrow V(s)$

\qquad $V(s)\leftarrow \max_a\sum_{s',r}p(s',r|s,a)\left[r+\gamma V(s')\right]$

\qquad $\Delta\leftarrow \max(\Delta,|v-V(s)|)$

\quad 直到 $\Delta<\theta$

\medskip
\noindent 输出一个确定性策略 $\pi\approx\pi_*$，满足
\[
\pi(s)=\operatorname*{arg\,max}_a\sum_{s',r}p(s',r|s,a)\left[r+\gamma V(s')\right].
\]
\end{minipage}}
\end{center}

价值迭代实际上在每一次扫掠中结合了一次策略评估扫掠和一次策略改进扫掠。
通常，在每次策略改进扫掠之间插入多次策略评估扫掠，可以获得更快的收敛。
一般来说，整个截断策略迭代算法类都可以看成一列扫掠，其中一些使用策略评估更新，另一些使用价值迭代更新。
因为 (4.10) 中的最大化操作是这些更新之间唯一的差别，这只是意味着最大化操作被加入到了若干策略评估扫掠中。
所有这些算法都会在折扣有限 MDP 上收敛到最优策略。

\noindent\textbf{例 4.3：赌徒问题}
一个赌徒有机会对一系列抛硬币的结果下注。
如果硬币正面朝上，他赢得与本次下注额相同的美元数；如果反面朝上，他输掉下注额。
当赌徒通过达到 $100$ 美元目标而获胜，或者因没钱而失败时，游戏结束。
在每次抛硬币时，赌徒必须决定拿出资本中的哪一部分下注，单位为整数美元。
这个问题可以表述为一个无折扣、情节式、有限 MDP。
状态是赌徒的资本，$s\in\{1,2,\ldots,99\}$；动作是下注额，$a\in\{0,1,\ldots,\min(s,100-s)\}$。
除赌徒达到目标的转移之外，所有转移的奖励都是零；达到目标时奖励为 $+1$。
状态价值函数给出从每个状态获胜的概率。
策略是从资本水平到下注额的映射。
最优策略最大化达到目标的概率。
令 $p_h$ 表示硬币正面朝上的概率。
如果 $p_h$ 已知，那么整个问题就是已知的，并且可以例如用价值迭代来求解。
图 4.3 展示了当 $p_h=0.4$ 时，价值函数在价值迭代的连续扫掠中的变化，以及最终找到的策略。
这个策略是最优的，但并不唯一。
事实上，对于相对于最优价值函数的 $\operatorname*{arg\,max}$ 动作选择，有一整族对应的最优策略。
你能猜出整个族是什么样的吗？

\begin{figure}[htbp]
\centering
\includegraphics[width=0.62\textwidth]{tex/强化学习/figures/rl5-fig-4-3.png}
\caption{当 $p_h=0.4$ 时赌徒问题的解。上图显示了由价值迭代的连续扫掠找到的价值函数。下图显示了最终策略。}
\end{figure}

\noindent\textbf{习题 4.8}
为什么赌徒问题的最优策略具有这样奇特的形式？
特别是，当资本为 $50$ 时，它一次性全押；但当资本为 $51$ 时，它却不这样做。
为什么这是一个好策略？

\noindent\textbf{习题 4.9（编程）}
实现赌徒问题的价值迭代，并分别在 $p_h=0.25$ 和 $p_h=0.55$ 时求解。
在编程时，你可能会发现引入两个哑状态很方便，它们对应于资本为 $0$ 和 $100$ 的终止，价值分别为 $0$ 和 $1$。
请像图 4.3 那样，以图形方式展示你的结果。
当 $\theta\to0$ 时，你的结果稳定吗？

\noindent\textbf{习题 4.10}
动作价值 $q_{k+1}(s,a)$ 的价值迭代更新 (4.10) 的类似形式是什么？

\subsection*{4.5\quad 异步动态规划}

到目前为止，我们讨论过的 DP 方法有一个主要缺点：它们涉及对 MDP 的整个状态集合进行操作，也就是说，它们需要对状态集合进行扫掠。
如果状态集合很大，那么即使一次扫掠也可能昂贵到不可承受。
例如，西洋双陆棋有超过 $10^{20}$ 个状态。
即使我们每秒能在一百万个状态上执行一次价值迭代更新，完成一次扫掠也需要一千多年。

\textbf{异步 DP} 算法是原地迭代 DP 算法，但它们并不按照状态集合的系统性扫掠来组织。
这些算法以任意顺序更新状态的价值，使用当时可用的其他状态价值。
一些状态的价值可能在另一些状态的价值更新一次之前就被更新多次。
不过，为了正确收敛，异步算法必须持续更新所有状态的价值：它不能在计算中的某个时刻之后忽略任何状态。
异步 DP 算法在选择要更新的状态方面提供了很大的灵活性。

例如，一种异步价值迭代版本在每个步骤 $k$ 只原地更新一个状态 $s_k$ 的价值，使用价值迭代更新 (4.10)。
如果 $0\le\gamma<1$，只要所有状态都在序列 $\{s_k\}$ 中出现无限多次（该序列甚至可以是随机的），就能保证渐近收敛到 $v_*$。
在无折扣的情节式情形中，可能存在某些更新顺序不会导致收敛，但这些顺序相对容易避免。
类似地，可以把策略评估更新和价值迭代更新混合起来，以产生一种异步截断策略迭代。
虽然这个算法以及其他更不常见的 DP 算法细节超出了本书范围，但可以清楚看出，少数不同的更新构成了构件，它们可以在各种无扫掠 DP 算法中灵活使用。

当然，避免扫掠并不必然意味着计算量更少。
它只是意味着算法不需要陷入某次漫长扫掠，直到扫掠结束才能在改进策略上取得进展。
我们可以利用这种灵活性来选择应用更新的状态，从而提高算法的推进速度。
我们可以尝试安排更新顺序，使价值信息以高效方式从一个状态传播到另一个状态。
有些状态可能不需要像其他状态那样频繁地更新价值。
如果某些状态与最优行为无关，我们甚至可以尝试完全跳过它们。
第 8 章讨论了实现这些想法的一些方式。

异步算法也使得把计算与实时交互混合起来更容易。
为了解决一个给定 MDP，我们可以在智能体实际经历该 MDP 的同时运行一个迭代 DP 算法。
智能体的经验可以用来决定 DP 算法对哪些状态应用更新。
与此同时，DP 算法给出的最新价值和策略信息可以引导智能体的决策。
例如，可以在智能体访问某些状态时对这些状态应用更新。
这使得可以把 DP 算法的更新集中到状态集合中与智能体最相关的部分。
这种聚焦是强化学习中反复出现的主题。

\subsection*{4.6\quad 广义策略迭代}

策略迭代由两个同时存在、相互作用的过程组成：一个过程使价值函数与当前策略一致（策略评估），另一个过程使策略相对于当前价值函数变得贪心（策略改进）。
在策略迭代中，这两个过程交替进行，每个过程在另一个过程开始之前完成；但这并不是真正必要的。
例如，在价值迭代中，每次策略改进之间只执行一次策略评估迭代。
在异步 DP 方法中，评估过程和改进过程以更细的粒度交错。
在某些情形中，一个状态会在返回另一个过程之前在一个过程中被更新。
只要两个过程都持续更新所有状态，最终结果通常是相同的：收敛到最优价值函数和最优策略。

我们使用\textbf{广义策略迭代}（generalized policy iteration, GPI）这个术语，来指让策略评估过程和策略改进过程相互作用的一般思想，而不依赖于两个过程的粒度和其他细节。
几乎所有强化学习方法都可以很好地描述为 GPI。
也就是说，它们都有可识别的策略和价值函数，策略总是相对于价值函数被改进，而价值函数总是被推向该策略的价值函数，正如下图所示。

\begin{figure}[htbp]
\centering
\includegraphics[width=0.28\textwidth]{tex/强化学习/figures/rl5-gpi-cycle.png}
\caption*{广义策略迭代中评估与改进两个过程的相互作用。}
\end{figure}

如果评估过程和改进过程都稳定下来，也就是说，不再产生变化，那么价值函数和策略必定是最优的。
价值函数只有在与当前策略一致时才会稳定，而策略只有在相对于当前价值函数贪心时才会稳定。
因此，两个过程只有在找到一个相对于其自身评估函数为贪心的策略时才会同时稳定。
这意味着贝尔曼最优方程 (4.1) 成立，因此策略和价值函数都是最优的。

GPI 中的评估过程和改进过程可以看作既相互竞争又相互合作。
它们相互竞争，是因为它们把系统拉向相反方向。
使策略相对于价值函数贪心，通常会让价值函数对于改变后的策略而言变得不正确；而使价值函数与策略一致，通常又会使该策略不再是贪心的。
不过从长远看，这两个过程相互作用，以找到一个单一的联合解：最优价值函数和一个最优策略。

也可以把 GPI 中评估过程和改进过程的相互作用理解为两个约束或目标之间的作用。
例如，可以把它们看成二维空间中的两条直线，如下图所示。
虽然真实几何关系要复杂得多，但这个图暗示了真实情形中发生的事情。
每个过程都把价值函数或策略推向表示两个目标之一的某条直线。
两个目标会相互作用，因为这两条直线并不正交。
直接朝一个目标推进会导致远离另一个目标的一些移动。
然而，不可避免地，联合过程会更接近最优性这一总体目标。
图中的箭头对应于策略迭代的行为：每个箭头都把系统一路带到完全达到两个目标之一。
在 GPI 中，也可以朝每个目标采取更小的、不完整的步骤。
无论哪种情形，两个过程共同实现了最优性的总体目标，尽管其中任何一个过程都不是在直接试图实现它。

\begin{figure}[htbp]
\centering
\includegraphics[width=0.44\textwidth]{tex/强化学习/figures/rl5-gpi-constraints.png}
\caption*{把 GPI 看成两个相互作用的目标：价值函数与策略一致，以及策略相对于价值函数贪心。}
\end{figure}

\subsection*{4.7\quad 动态规划的效率}

DP 对非常大的问题可能并不实用，但与其他求解 MDP 的方法相比，DP 方法实际上相当高效。
如果忽略若干技术细节，那么 DP 方法找到一个最优策略所需的（最坏情形）时间，是状态数和动作数的多项式函数。
如果 $n$ 和 $k$ 分别表示状态数和动作数，这意味着 DP 方法所需的计算操作数小于 $n$ 和 $k$ 的某个多项式函数。
虽然（确定性）策略的总数是 $k^n$，但 DP 方法保证能在多项式时间内找到最优策略。
在这个意义上，DP 比任何在策略空间中直接搜索的方法都指数级地更快，因为直接搜索必须穷尽地检查每一个策略，才能提供同样的保证。
线性规划方法也可以用来求解 MDP，并且在某些情况下，它们的最坏情形收敛保证优于 DP 方法。
但是，在线性规划方法变得不实际时，状态数量要比 DP 方法小得多（大约差一个 $100$ 倍的因子）。
对于最大的问题，只有 DP 方法是可行的。

有时人们认为 DP 的适用性有限，是因为所谓\textbf{维数灾难}：状态数通常随状态变量的个数指数增长。
大的状态集合确实会造成困难，但这是问题本身固有的困难，而不是 DP 作为求解方法的困难。
事实上，与直接搜索和线性规划等竞争方法相比，DP 更适合处理大的状态空间。

在实践中，借助今天的计算机，DP 方法可以用来求解具有数百万个状态的 MDP。
策略迭代和价值迭代都被广泛使用，而且一般来说尚不清楚二者哪一个更好。
实践中，这些方法通常比它们理论上的最坏情形运行时间收敛得快得多，特别是当它们从较好的初始价值函数或策略开始时。

对于具有大状态空间的问题，通常更偏好异步 DP 方法。
即使完成一次同步方法的扫掠，也需要为每个状态进行计算并存储信息。
对某些问题来说，即使这种程度的内存和计算也并不实际；然而问题仍然可能是可解的，因为在最优解轨迹上出现的状态相对较少。
异步方法和其他 GPI 变体可以应用于这类情形，并且可能比同步方法更快地找到好的或最优的策略。

\subsection*{4.8\quad 小结}

在本章中，我们熟悉了动态规划中与求解有限 MDP 相关的基本思想和算法。
\textbf{策略评估}指对给定策略的价值函数进行（通常是）迭代计算。
\textbf{策略改进}指在给定某策略价值函数的情况下计算一个改进后的策略。
把这两种计算合在一起，我们得到策略迭代和价值迭代，这是两种最流行的 DP 方法。
在给定 MDP 完整知识的情况下，二者都可以可靠地计算有限 MDP 的最优策略和价值函数。

经典 DP 方法以扫掠方式遍历状态集合，对每个状态执行一次期望更新操作。
每个这样的操作都根据所有可能后继状态的价值及其发生概率来更新一个状态的价值。
期望更新与贝尔曼方程密切相关：它们不过是把这些方程转化成赋值语句。
当更新不再造成任何价值变化时，就已经收敛到满足相应贝尔曼方程的价值。
正如有四个主要价值函数（$v_\pi$、$v_*$、$q_\pi$ 和 $q_*$）一样，也有四个相应的贝尔曼方程和四个相应的期望更新。
对 DP 更新操作的直观理解由它们的备份图给出。

把 DP 方法以及事实上几乎所有强化学习方法看成广义策略迭代（GPI），可以获得深刻理解。
GPI 是两个围绕近似策略和近似价值函数相互作用的过程的一般思想。
一个过程把策略视为给定并执行某种形式的策略评估，把价值函数改得更像该策略的真实价值函数。
另一个过程把价值函数视为给定并执行某种形式的策略改进，在假定价值函数就是其价值函数的情况下把策略改得更好。
虽然每个过程都会改变另一个过程的基础，但总体上它们共同寻找一个联合解：一个策略和价值函数，它们不被任何一个过程改变，并因此是最优的。
在一些情形中，可以证明 GPI 收敛，最显著的是本章介绍的经典 DP 方法。
在其他情形中，收敛性尚未得到证明，但 GPI 的思想仍然能改进我们对这些方法的理解。

并不一定要以完整扫掠方式穿过状态集合来执行 DP 方法。
异步 DP 方法是原地迭代方法，它们以任意顺序更新状态，这个顺序也许由随机方式决定，并且可能使用过时的信息。
许多这样的方法可以看作 GPI 的细粒度形式。

最后，我们注意到 DP 方法还有一个特殊性质。
所有 DP 方法都根据对后继状态价值的估计来更新状态价值的估计。
也就是说，它们基于其他估计来更新估计。
我们把这个一般思想称为\textbf{自举}（bootstrapping）。
许多强化学习方法都会执行自举，即使它们不像 DP 那样需要环境的完整而准确的模型。
在下一章中，我们考察不需要模型且不自举的强化学习方法。
再下一章中，我们考察不需要模型但会自举的方法。
这些关键特征和性质可以分离开来，但也可以以有趣的方式混合在一起。

\section*{文献与历史备注}

“动态规划”这个术语归功于 Bellman（1957a），他展示了这些方法如何应用于范围广泛的问题。
许多教材都对 DP 作了详尽论述，包括 Bertsekas（2005, 2012）、Bertsekas 和 Tsitsiklis（1996）、Dreyfus 和 Law（1977）、Ross（1983）、White（1969）以及 Whittle（1982, 1983）。
我们对 DP 的兴趣仅限于它在求解 MDP 中的用途，但 DP 也适用于其他类型的问题。
Kumar 和 Kanal（1988）对 DP 作了更一般的考察。

据我们所知，DP 与强化学习之间的第一次联系，是 Minsky（1961）在评论 Samuel 的跳棋程序时建立的。
在一个脚注中，Minsky 提到，可以把 DP 应用于 Samuel 的备份过程能够用封闭解析形式处理的问题。
这句话可能误导了人工智能研究者，使他们相信 DP 只限于解析上可处理的问题，因此与人工智能基本无关。
Andreae（1969b）在强化学习的语境中提到了 DP，特别是策略迭代，尽管他没有具体建立 DP 与学习算法之间的联系。
Werbos（1977）提出了一种近似 DP 的方法，称为“启发式动态规划”，它强调用于连续状态问题的梯度下降方法（Werbos, 1982, 1987, 1988, 1989, 1992）。
这些方法与本书讨论的强化学习算法密切相关。
Watkins（1989）明确把强化学习与 DP 联系起来，把一类强化学习方法刻画为“增量动态规划”。

\noindent\textbf{4.1--4}
这些节描述的是成熟的 DP 算法，前面引用的任何一般 DP 参考文献都会覆盖它们。
策略改进定理和策略迭代算法归功于 Bellman（1957a）和 Howard（1960）。
我们的表述受到 Watkins（1989）所采用的策略改进局部视角影响。
我们把价值迭代讨论为截断策略迭代的一种形式，是基于 Puterman 和 Shin（1978）的方法；他们提出了一类称为\textbf{修正策略迭代}的算法，其中策略迭代和价值迭代都是特殊情形。
Bertsekas（1987）给出了说明如何使价值迭代在有限时间内找到最优策略的分析。

迭代策略评估是求解线性方程组的经典逐次逼近算法的一个例子。
使用两个数组的算法版本，其中一个数组保存旧值而另一个数组被更新，通常称为 Jacobi 式算法，这是沿用 Jacobi 对该方法的经典使用。
它有时也称为同步算法，因为它的效果就像所有价值都同时更新一样。
第二个数组用于顺序地模拟这种并行计算。
算法的原地版本通常称为 Gauss--Seidel 式算法，这是沿用用于求解线性方程组的经典 Gauss--Seidel 算法。
除了迭代策略评估之外，其他 DP 算法也可以实现为这些不同版本。
Bertsekas 和 Tsitsiklis（1989）对这些变体及其性能差异作了出色论述。

\noindent\textbf{4.5}
异步 DP 算法归功于 Bertsekas（1982, 1983），他也称它们为分布式 DP 算法。
异步 DP 的最初动机，是在多处理器系统上实现 DP；处理器之间存在通信延迟，并且没有全局同步时钟。
Bertsekas 和 Tsitsiklis（1989）对这些算法作了广泛讨论。
Jacobi 式和 Gauss--Seidel 式 DP 算法是异步版本的特殊情形。
Williams 和 Baird（1990）提出的 DP 算法在比我们讨论的粒度更细的层次上是异步的：更新操作本身被拆分成可以异步执行的步骤。

\noindent\textbf{4.7}
本节在 Michael Littman 的帮助下写成，并基于 Littman、Dean 和 Kaelbling（1995）。
“维数灾难”这个短语归功于 Bellman（1957a）。
强化学习中线性规划方法的基础工作由 Daniela de Farias 完成（de Farias, 2002; de Farias and Van Roy, 2003）。

\end{document}
